% Included by notes/note.tex; compile the integrated note from the project root.
\clearpage
\part{Circuit-level fixed-fiber theory}
\label{part:circuit}
\section{Circuit faults, detector maps, and observable semantics}
\label{sec:cl-dem}
This part concerns a finite, fully terminated triangular 6.6.6 memory
experiment, one chosen color, and the sector detected by $Z$ checks and
read out as logical $Z$. Its mathematical input is the actual retained
labelled stage-2 DEM. The graph results below also hold for any finite
one-observable graphlike DEM satisfying their explicit hypotheses. They
are not a claim about an open future boundary or sliding-window decoding.
All new results are our derivations or explicit adaptations, not priority
claims. The earlier perfect-measurement results remain unchanged.

\begin{definition}[DEM algebra]\label{def:cl-dem}
Flatten a finite Stim DEM, resolving detector offsets and repetitions.
Give each error instruction a stable ID $e\in E_0$ and retain its probability
$q_e$, detector parity vector $D_e$, and observable parity vector $O_e$.
Set $B[:,e]=D_e$, $L[:,e]=O_e$. Then
\[
 B:\F^{E_0}\longrightarrow\F^{V_{\rm det}},\qquad
 L:\F^{E_0}\longrightarrow\F^{V_{\rm obs}},\qquad
 (s,\ell)=(Bx,Lx).
\]
An instruction containing separator symbols describes one correlated event;
its components' targets are XORed to form this column. Separators do not
supply independent Bernoulli variables. Detector coordinates are metadata,
not the definition of either matrix. Observable IDs are frame-change
labels, not additional constrained detector rows.
\end{definition}
These are Stim's error-instruction semantics \cite{StimDEM}.
For a fixed Clifford circuit with Pauli faults, let $A_{\rm rec}$ send
faults to changes of measurement-record bits, and let $P_{\rm det}$ and
$P_{\rm obs}$ select the annotated record parities. With ideal reference
signs subtracted,
\begin{equation}
 B=P_{\rm det}A_{\rm rec},\qquad L=P_{\rm obs}A_{\rm rec}.
 \label{eq:cl-record}
\end{equation}
This follows by composing binary parity maps. It is the detector-matrix
factorization of Derks et al.\ \cite[Definitions 2.6, 2.8--2.9]{DerksDEM},
with separate symbols for record checks and fault incidence. A DEM may
already compress physically distinct faults; its columns are not necessarily
individual gate locations. Retain a relation to circuit locations, channel
outcomes, and separator groups separately. General disjoint Pauli channels
need not admit an exact independent DEM. Record Stim's conversion options
and any approximation, and never infer exact probabilities from the binary
identities alone.

For the memory in Lee, physical detector rows are the successive check
parities of Eq.\ (6), including the initial $Z$ checks and final ancilla/data
consistency checks; the observable is the final side parity of Eq.\ (7)
\cite[Sec.~4.1]{Lee}. The $X$-check sector lacks those same initial/final
rows in a $Z$ memory. Detector Pauli type names the measured check, not
the error Pauli that anticommutes with it.

\section{Lee decomposition, compression, and the actual implementation}
\label{sec:cl-decomp}
\subsection{Paper algorithm with typed metadata}
Lee's Algorithm~1 first separates each error into its $Z$-detector and
$X$-detector parts, sending each observable to its corresponding measured
Pauli sector. Each retained part keeps the original probability. The
paper's line~6 retains a part only if it has a nonempty detector set.
Thus a pure observable-only part is omitted by the literal algorithm;
its absence must be checked, or its loss recorded. Line~5's phrase
``observables in $\mathcal D$'' is read using $\mathcal O$, consistently
with the input tuple and the prose following the algorithm.
Compression at line~10 uses the \emph{complete} target set, including
observable IDs. For independent members of one such class $C$ it gives
\begin{equation}
 q_C=\frac{1-\prod_{e\in C}(1-2q_e)}2.
 \label{eq:cl-compression}
\end{equation}
To derive this, the expectation of $(-1)^{\oplus_{e\in C}x_e}$ is the
product of the expectations $1-2q_e$. This proves exact parity compression
within that independent model, including the displayed two-event formula.
It is not a minimum-cost representative operation.

For a compressed separated mechanism $e$, write $A_e$ for its physical
$c$ detector targets, $N_e$ for its non-$c$ targets, and $o_e$ for its
observable vector. In the paper the restricted model first collects every
$e$ with $0<|N_e|\le2$, erases its observable labels, then compresses by
$N_e$. Give each distinct nonempty retained target set $N$ a restricted
edge ID and a constrained virtual row $v_N$. Let $H_{\rm rest}[:,N]=N$.
The restricted observable erasure is intentional: stage~1 predicts edge
occupancy, not a physical logical observable.

The stage-2 rule retains precisely the following mechanisms from the
separated input:
\begin{equation}
 d_e=\begin{cases}
 A_e,&N_e=\varnothing,\ |A_e|\le2,\\
 A_e+v_{N_e},&0<|N_e|\le2,\ |A_e|\le1.
 \end{cases}
 \qquad \ell_e=o_e .                         \label{eq:cl-decomp}
\end{equation}
All others are removed. Each retained stage-2 column keeps its own
probability; it does not inherit the compressed restricted probability.
There is no instruction to forget observable labels when two stage-2
edges have the same endpoints.

\begin{definition}[Stage-2 maps and row types]\label{def:cl-stage2}
Let $E_c^{\rm circ}$ be the retained mechanism IDs and
$V_c^{\rm real}=V_c^{\rm phys}\sqcup V_c^{\rm virt}$, a tagged union.
Physical rows are measured $c$ detectors; virtual rows are the stage-1
predicted restricted-edge parities. Both are \emph{constrained} rows in
stage~2. Put
\[
 D_c^{\rm circ}[:,e]=d_e,\qquad L_c^{\rm circ}[:,e]=\ell_e,
 \qquad \omega_c(e)=\log\frac{1-q_e}{q_e}.
\]
The metric uses supplied finite nonnegative weights, hence ordinarily
$0<q_e\le1/2$. Zero-probability mechanisms are excluded from the admissible
fault set; negative log odds are outside this metric theorem. This is
not a prescription to clip negative hard-decoder weights.
Inactive zero rows used for storage padding are recorded separately and
may be omitted in analysis only if their supplied syndrome is zero.
The artificial matching boundary has no row in $D_c^{\rm circ}$.
Below $D,L_2,w$ abbreviate these maps and weights locally.
\end{definition}

Define $Q_{\rm virt}[:,e]=v_{N_e}$ if $N_e\ne\varnothing$, and zero
otherwise. If $B_{\rm ret}$ is the physical detector matrix on retained
separated mechanisms, then columnwise
\begin{equation}
 D=\begin{pmatrix}B_c\\Q_{\rm virt}\end{pmatrix},\qquad
 B_{\rm ret}=\Lambda_c^{\rm circ}D,\qquad
 \Lambda_c^{\rm circ}(a,v)=(a,H_{\rm rest}v).
 \label{eq:cl-factor}
\end{equation}
The physical rows in the last expression are ordered as $c$, then non-$c$.
This proof uses $H_{\rm rest}v_{N_e}=N_e$ and is valid on every retained
column. Consequently a feasible stage~1 result $t$ and feasible stage~2
result $x$ with $Dx=(s_c^{\rm phys},t)$ reconstruct the supplied physical
syndrome through $\Lambda_c^{\rm circ}$. Existence of both feasible outputs
is a separate condition; filtering can destroy it.

\begin{proposition}[Observable-faithfulness and its limits]
\label{prop:cl-faithful}
Within the retained separated DEM, the observable of any mechanism vector
$x$ is exactly $L_2x$. Full-target compression preserves detector and
observable parity under its XOR projection. Graphlike filtering preserves
the labels of surviving columns, but need not preserve the logical quotient,
the minimum logical cost, or the noise distribution of the input model.
\end{proposition}
\begin{proof}
Every stage-2 substitution in \eqref{eq:cl-decomp} leaves $o_e$ unchanged;
linearity proves the first assertion. Let $P_{\rm comp}x$ XOR all original
variables in each full-target class. Equal columns give
$B_{\rm old}=B_{\rm comp}P_{\rm comp}$ and
$L_{\rm old}=L_{\rm comp}P_{\rm comp}$. Equation~\eqref{eq:cl-compression}
is exact for independent classes. Filtering instead restricts the domain
to a coordinate subspace. It can remove every odd kernel vector: take two
columns both incident on three detectors, with labels zero and one. Their
sum is invisible and odd, while deleting both leaves no logical class.
\end{proof}
The many-to-one XOR projection has no canonical inverse. A provenance
membership vector is not a physical fault lift: choosing \emph{all} members
of a two-element compressed class yields even, not odd, target parity.
A physical interpretation of a selected compressed variable requires a
choice of a parity-odd preimage, with its own probability and cost.
X/Z separation replaces a common event by independent effective events;
it preserves sector target labels but loses intersector correlations.
Neither that operation nor conditioning on a predicted stage~1 result
establishes physical statistical independence.

\subsection{Pinned source realization and concrete discrepancies}
The inspection here is at color-code-stim commit
\texttt{ba6f7dc} (full hash in the bibliography), not the earlier Phase-1
snapshot \cite{CircuitImplementation}. For ordinary memory,
\texttt{DemManager.\_generate\_dem} first calls
\texttt{separate\_depolarizing\_errors} on the \emph{circuit}, then extracts
its flattened DEM. It does not first extract the correlated original DEM.
The one-qubit replacement has sector error probability $2p/3$.
For two-qubit depolarization each of $X_A,X_B,X_AX_B$ (and separately the
three $Z$ mechanisms) has probability
$(1-\sqrt{1-16p/15})/2$. These independent sector generators reproduce the
sector marginal: each nonzero two-bit outcome has probability $4p/15$,
since $a(1-a)=4p/15$ for $a=(1-\sqrt{1-16p/15})/2$.
The joint X/Z distribution changes. The implementation
also drops components at or below its $10^{-15}$ probability cutoff.
These details must be part of the recorded effective-model definition.

In \texttt{DemDecomp.decompose\_org\_dem}, the dictionary key contains
both detector IDs and observable strings. For already separated, unique
full-target inputs, each retained stage-2 mechanism has one source in the
manager's effective DEM; probability sorting permutes its source-map rows
consistently. Then $L_2=L_{\rm manager}J$ and
$B_{\rm manager}J=\Lambda_c^{\rm circ}D$, where $J$ is that column selection.
The outer decoder computes its observable using this map and the manager's
observable matrix. This establishes agreement with the hard-decoder frame
for that domain; the manager's ``original'' DEM is already noise-modified.

Two additional limits follow directly from the code. First, its
mixed-sector fallback assigns observables to the $Z$ part and assigns a
fresh singleton source list to each target key, overwriting an earlier
identical key rather than parity-compressing both sources. The later
stage-2 map rejects multiple rows referring to one original source.
Thus this fallback is not a general implementation of Lee's compression
or multi-observable Pauli separation. Second, its joint filtering happens
\emph{before} restricted probabilities are accumulated: a mechanism with
$|N_e|\le2$ but $|A_e|\ge2$ contributes to paper stage~1 but not to code
stage~1. The two algorithms can therefore select different fibers even
though \eqref{eq:cl-factor} holds for the retained code model.
No external source is changed by this theory task. A future adapter must
consume the actual frozen matrices and check the two displayed identities,
rather than silently replacing the hard decoder with the paper algorithm.

\section{The fixed-fiber logical quotient and existence of two classes}
\label{sec:cl-quotient}
\begin{theorem}[Fixed-fiber logical classification]\label{thm:cl-quotient}
Fix stage~1 and a nonempty stage-2 fiber $\mathcal F_s=\{x:Dx=s\}$.
Define
\[
 K_c^{\rm circ}=\ker D,\qquad R_c^{\rm circ}=\ker D\cap\ker L_2,
 \qquad \mathcal Q_c^{\rm circ}=K_c^{\rm circ}/R_c^{\rm circ}.
\]
Then $[z]\mapsto L_2z$ is a canonical isomorphism
\begin{equation}
 \mathcal Q_c^{\rm circ}\cong\operatorname{im}(L_2|_{\ker D}).
 \label{eq:cl-quotient}
\end{equation}
For one observable the image is $\F$ if and only if
\begin{equation}
 L_2\notin\operatorname{rowspan}D,
 \quad\text{equivalently}\quad
 \operatorname{rank}\begin{pmatrix}D\\L_2\end{pmatrix}
       =\operatorname{rank}D+1.                       \label{eq:cl-rank}
\end{equation}
In that case every nonempty fiber has exactly two frame classes.
\end{theorem}
\begin{proof}
The restricted logical map has kernel exactly $R_c^{\rm circ}$, so its
quotient map is well-defined, injective, and onto its stated image.
The annihilator of $\ker D$ is $\operatorname{rowspan}D$: inclusion is
immediate and both dimensions equal $\operatorname{rank}D$. A binary
one-row functional is either zero or surjective on a subspace, proving
the criterion. Translation by any $f\in\mathcal F_s$ identifies that
fiber with $f+\ker D$. Two candidates differ invisibly under $D$ and have
same or opposite frame according to $L_2(x+x')=0$ or $1$.
\end{proof}
The denominator means no observable-frame change within this effective
model. It is not asserted to be the group of static face stabilizers,
nor the set of all physically harmless circuit faults for unmeasured
logical observables. The other sector in a one-$Z$-observable memory
can have the zero logical map.

\begin{proposition}[A sufficient memory-circuit witness condition]
\label{prop:cl-memory}
Suppose the retained measured-sector model contains the effects of every
single data-qubit readout flip at the final boundary, with nonzero
probability, final face detectors as in Lee Eq.\ (6a), and logical parity
as in Eq.\ (7). Full-target compression may replace an effect by an
identical effect. Then \eqref{eq:cl-rank} holds for every color and every
odd $d\ge3$, for any number of rounds and any extraction schedule satisfying
these final detector definitions.
\end{proposition}
\begin{proof}
A readout flip at $q$ changes exactly the final real faces containing $q$
and the final logical side parity if that side contains $q$.
Its non-$c$ face set identifies its incident primal $c$ edge, except at
the opposite corner where that set is empty. Thus substitution gives
exactly the Phase-1 column $(H_cq,M_cq)$ on final physical and final
virtual rows. These columns pass the graphlike filter by
Lemma~\ref{lem:incidence}. Distinct qubits give distinct such columns:
the two endpoints of a primal $c$ edge have different surviving $c$-face
incidence (a face of color $c$ does not contain that edge); the unique
missing-edge corner is separate. Hence compression does not identify
these qubits with each other. Take the Phase-1 path supported on either
non-$c$ physical side. Its virtual and physical incidences cancel, and
its final logical parity is one by Lemma~\ref{lem:sideparity}. The same
binary combination of retained effects is an odd element of $\ker D$.
\end{proof}
This is a sufficient all-distance condition, not an assumption that every
noise model supplies readout faults. For absent or removed mechanisms use
the necessary-and-sufficient rank gate; do not infer two classes merely
from the number of encoded qubits. No noisy-circuit threshold enters this
proof.

\section{Logical cochains and the precise spacetime interpretation}
\label{sec:cl-geometry}
\begin{definition}[Algebraic complex and correlation cochain]
\label{def:cl-complex}
Choose a basis of $R_c^{\rm circ}$ and let $A_R$ be its inclusion into
$\F^{E_c^{\rm circ}}$. The complex is
\[
 U_R\xrightarrow{A_R}\F^{E_c^{\rm circ}}\xrightarrow{D}\F^{V_c^{\rm real}},
 \qquad DA_R=0.
\]
Its cochain maps are $D^{\mathsf T}$ and $A_R^{\mathsf T}$. The cochain
$\lambda_c=L_2^{\mathsf T}$ satisfies $A_R^{\mathsf T}\lambda_c=0$.
It pairs with $z\in\ker D$ by $\lambda_c^{\mathsf T}z=L_2z$ and
therefore defines a class in $H^1$ dual to the quotient in
\eqref{eq:cl-quotient}. Equivalent representatives are
$\lambda'_c=\lambda_c+D^{\mathsf T}g$.
\end{definition}
This complex explicitly quotients out frame-trivial differences; it does
not prove that the basis vectors of $U_R$ are local geometric two-cells.
For the bare graph complex without two-cells every edge cochain is a
cocycle, but its $H_1$ is the larger space $\ker D$.

There is also a circuit-level correlation pairing before decomposition.
Pull the logical record selector backward through the circuit:
$\Sigma_L^{\rm fault}=A_{\rm rec}^{\mathsf T}P_{\rm obs}^{\mathsf T}$.
For every fault chain $x$,
\[
 \langle x,\Sigma_L^{\rm fault}\rangle
 =P_{\rm obs}A_{\rm rec}x=Lx.                 \label{eq:cl-pairing}
\]
This is an explicit cochain pairing, derived from \eqref{eq:cl-record}.
Clifford propagation realizes it as the symplectic commutation pairing
with backward-propagated measurement operators; forward/backward
adjointness is also proved in Delfosse--Paetznick
\cite[Proposition~3, Sec.~3.4]{SpacetimeCodes}.
If an effective fault lift $J_{\rm fault}$ into a specified physical fault
complex is supplied with $L_{\rm fault}J_{\rm fault}=L_2$ and the matching
detector identity, then pullback gives
$\langle J_{\rm fault}z,\Sigma_L^{\rm fault}\rangle=L_2z$.
A compression membership relation alone does not meet this hypothesis.
A geometric surface is a realization of this cochain only after its
intersection functional agrees on all columns, or differs by $g^{\mathsf T}D$
on kernel chains. This states an executable parity check for a proposed
surface, not a proof from its drawing.

\subsection{What the literature does and does not identify}
Hillmann et al.\ define fault complexes and express CSS foliation using
a repetition-code tensor product \cite[Eqs.~(4)--(6), p.~2]{FaultComplex}.
Logical correlations and undetected faults are paired homology/cohomology
classes in that construction. This is existing theory for its specified
complex, not a theorem that Lee's compressed virtual rows are local cells.
Delfosse--Paetznick construct spacetime stabilizer codes for Clifford
circuits \cite[Theorems~1--3]{SpacetimeCodes}; their maximum-likelihood
correspondence uses its corresponding fault model, not Lee's independently
reweighted model. Bomb\'in et al.\ distinguish resource and syndrome
complexes; their color-code fusion complex has three-colorable vertices
and triangular two-cells \cite[Appendix~F, Definition~7]{FaultTolerant}.
Neither construction identifies an artificial matching boundary with a
physical boundary.

Herzog et al.\ explicitly represent a triangular 6.6.6 memory as a prism,
with temporal preparation/readout faces and vertical logical correlation
surfaces \cite[Sec.~III.A--B, Figs.~5--6]{ColorPipes}. Their standard prism
uses $d$ rounds; changing the number of repetitions changes its time
extent. The prism is a protocol-level representation. It is not the
labelled $c$-only DEM, and its uncolored vertical walls are not Lee's
color-resolved matching terminals. These are source-backed geometric
interpretations; an exact local map from the actual compressed Lee model
to such a prism is \emph{open} here.

\begin{proposition}[Controlled product comparison]\label{prop:cl-product}
Let $C$ be the Phase-1 algebraic complex of Definition~\ref{def:complex}
and $I_T$ the cellular chain complex of an interval with vertices
$0,\ldots,T$ and edges $j=[j-1,j]$. The total tensor complex
$C\otimes I_T$ is chain-homotopy equivalent to $C$.
\end{proposition}
\begin{proof}
Let $p_T$ send every time vertex to $1$ and every time edge to zero;
let $i_T(1)=0$ (the basis vertex at time zero). Define
$h_T(j)=[0,1]+\cdots+[j-1,j]$ on time vertices, and zero on time edges.
Then $p_Ti_T=1$ and $\partial h_T+h_T\partial=1+i_Tp_T$, by cancellation
along the interval, including its two ends. Tensor these maps with the
identity on $C$. In characteristic two the cross terms in the homotopy
identity cancel. This proves the chain-homotopy equivalence.
\end{proof}
This is our direct algebraic comparison with the product intuition.
It supplies no cost isometry: different time layers can have different
weights. Gate faults, hooks, changed final constraints, filtering and
virtual substitutions require a further chain map; none is silently
identified with this ideal tensor product.

\subsection{Five boundary objects and temporal mechanisms}
Spatial color boundaries belong to the physical patch. Temporal boundaries
specify preparation and final record checks. The matching boundary is an
unconstrained algebraic endpoint. A correlation surface has its own
geometric boundary. An analysis cut introduces auxiliary terminal copies.
These five meanings are distinct. Virtual detectors have a restricted-edge
source set and a predicted bit, not a newly measured physical check.
Their coordinate metadata should store all source detector coordinates;
an average coordinate, if used for display, cannot assign their topology.

For example a $Z$-ancilla readout flip in round $1$ contributes to the
first and second $Z$ detectors (or to the first and final consistency
detectors if $T=1$), with observable label zero. A final data readout flip
contributes to its final incident face detectors and may flip the logical
side parity. It can become a half-edge after color decomposition.
The unmeasured sector and circuit variants can have additional end-time
half-edges. All are classified by $d_e,\ell_e$; their time coordinate or
surviving row type alone does not decide their logical action.
Closed memory therefore does not mean absence of temporal half-edges.
Removing future detector rows enlarges $\ker D$ and can change its logical
image; none of the memory conclusions authorizes sliding-window decoding.

\section{Fundamental residual minimum and the canonical logical cover}
\label{sec:cl-cover}
\begin{definition}[Boundary completion and residual minimum]
\label{def:cl-swim}
Complete each one-detector column at one artificial vertex $b_*$, retain
each two-detector column as an edge, and represent each zero-detector
column by a labelled loop at $b_*$. Call this multigraph $G_c^\bullet$.
No new edge connects distinct boundaries. Let $\bar\omega_c(e)\ge0$ be
the uncovered interval length of each original mechanism, specified in
Section~\ref{sec:cl-clusters}. Define, before selecting any cut,
\begin{equation}
 \boxed{\displaystyle \phi_c^{\rm circ}
 =\min_{Dz=0,\ L_2z=1}\sum_e\bar\omega_c(e)z_e.}
 \label{eq:cl-swim}
\end{equation}
An empty minimum is $+\infty$, accompanied by no logical witness.
\end{definition}
For any fixed $f$ with $Df=s$, the substitution $x=f+z$ proves that the
feasible set is exactly the detector-invisible opposite-frame fixed-fiber
difference set. Thus \eqref{eq:cl-swim} is its minimum contracted cost,
not necessarily the difference between uncontracted optimum weights.
It is independent of a drawing and depends on the chosen cluster data.

\begin{lemma}[Completion introduces no false logical chain]
\label{lem:cl-boundary}
A binary mechanism chain has $Dz=0$ if and only if its completed full
incidence vanishes in $G_c^\bullet$.
\end{lemma}
\begin{proof}
All real rows are identical. The sum of all incidence rows is zero,
including for loops, so zero real incidence forces zero incidence at
$b_*$. Conversely full zero incidence implies real zero incidence.
\end{proof}
In particular half-edges occur an even number of times in an invisible
chain. The completion is valid for this fixed matrix; a separate zero-cost
edge that changes sheet or carries a logical label is not permitted.

\begin{theorem}[Exact binary-cover algorithm]\label{thm:cl-cover}
Give edge $e$ the label $\lambda_e=L_2[e]$. The cover has vertices $(v,a)$,
$a\in\F$, and, for every base edge $e=uv$, two lifted edges
\[
 (u,a)\longleftrightarrow(v,a+\lambda_e),
 \qquad\text{each of cost }\bar\omega_c(e).
\]
Both lifts keep the source ID $e$. A loop of label one has two parallel
lifts between its two sheet copies; a zero-labelled loop stays a loop.
Then
\begin{equation}
 \phi_c^{\rm circ}
 =\min_{v\in V(G_c^\bullet)}
     \operatorname{dist}_{\widehat G_c}((v,0),(v,1)).
 \label{eq:cl-cover}
\end{equation}
A minimizing path projects, with traversals reduced modulo two, to a
witness $z$ with $Dz=0$, $L_2z=1$, and cost $\phi_c^{\rm circ}$.
For $n$ base vertices and $m$ labelled mechanisms this can be computed
in $O(n(m+n)\log(2+n))$ time and $O(m+n)$ working storage, given residuals.
\end{theorem}
\begin{proof}
Induction on a walk's length shows that its lift from $(v,a)$ ends on
sheet $a+\sum_e\lambda_e$ (counting multiplicity) over its endpoint.
Every path between opposite copies of the \emph{same} vertex therefore
projects to an odd closed walk. Its mod-two chain has zero incidence,
odd label and cost at most its walk cost by nonnegativity.

Conversely, an admissible chain is an even multigraph by
Lemma~\ref{lem:cl-boundary}. Pair its half-edges to decompose it into
edge-disjoint closed trails, then split at repeated vertices to obtain
cycles (allowing loops and parallel-edge two-cycles). The total label is
one, so some cycle is odd. Its cost is at most the chain's cost. Starting
at any vertex on this cycle lifts it to an opposite-sheet path or walk
of that cost. These two inequalities prove \eqref{eq:cl-cover}, including
zero costs and disconnected inputs, and force equality for the projected
minimizer. Dijkstra from each $(v,0)$ to $(v,1)$ on a $2n$-vertex,
$2m$-edge nonnegative graph gives the bound, reusing search storage.
Save source IDs in a predecessor tree; reduce their multiplicities XORwise.
\end{proof}
The construction is the standard voltage-cover idea; homology-signature
covers with inherited weights appear, for example, in
\cite[Sec.~5.1]{HomologyCovers}. The proof here needs just one label bit,
not a surface embedding or an identification of all homology classes.
For multiple observables a nonzero total label still has a nonzero cycle,
but a prescribed multi-bit class can require several cycles. No one-bit
complexity is asserted for that more general optimization.

One arbitrary source is insufficient: attach a unit-weight triangle with
odd total label to $b_*$ by a bridge of weight ten. The odd minimum is
three, while the distance from $(b_*,0)$ to $(b_*,1)$ is twenty-three.
A multisource search between \emph{all} zero and one sheet vertices is
also wrong: one labelled edge with two real endpoints can yield a finite
answer even when $\ker D=0$. Endpoints must project to the same base vertex.

\section{When a two-boundary cut is exact}
\label{sec:cl-cut}
A surface cut is not automatically a shortest path problem between arbitrary
points on its two sides. Those points may project to different constrained
vertices or interior edge points. Their projected path then has nonzero
boundary. A physical correlation surface must supply a valid chain-lifting
rule, not merely an intersection picture.

\begin{theorem}[Checkable two-terminal cut condition]\label{thm:cl-cut}
Let $G_{\rm int}$ contain all two-real-endpoint edges and real vertices of
$G_c^\bullet$. Suppose every cycle of $G_{\rm int}$ has label zero.
Equivalently there is a binary potential $g$ on real vertices with
$\lambda_{uv}=g(u)+g(v)$ on every internal edge. Extend $g(b_*)=0$ and
gauge labels by $\lambda'_e=\lambda_e+g^{\mathsf T}D[:,e]$.
Construct $G_c^{\rm cut}$ as follows: keep each internal edge once, replace
$b_*$ by two free terminals $b^0,b^1$, and attach a half-edge $e$ at
real vertex $u$ to $b^{\lambda'_e}$. Replace a zero-detector odd loop by
an edge $b^0b^1$, and an even loop by a loop at $b^0$. Keep all mechanism
IDs and residual costs. Then
\begin{equation}
 \phi_c^{\rm circ}
 =\operatorname{dist}_{\bar G_c^{\rm cut}}(b^0,b^1).
 \label{eq:cl-cut}
\end{equation}
Thus $B_c^0=\{b^0\}$ and $B_c^1=\{b^1\}$ are valid boundary sets.
The balance test, potential construction and shortest path take
$O((m+n)\log(2+n))$ time given residual costs.
\end{theorem}
\begin{proof}
Choose a spanning forest of $G_{\rm int}$ and assign $g$ by accumulated
labels from each root. A non-tree edge satisfies the equation precisely
when its fundamental cycle is even, proving the equivalence and giving
a linear-time check, including parallel edges. Internal gauged labels
are zero. On $\ker D$, gauging does not change logical parity.
The incidence at $b^1$ in the cut graph is exactly $L_2z$ on such chains;
handshake makes the incidence at $b^0$ the same. Odd chains therefore
have boundary $b^0+b^1$, and even chains have full boundary zero.
Every odd chain contains a simple terminal path by the odd-degree
component argument, and every such path is odd. Nonnegative costs prove
the distance equality. Original edge incidence is otherwise unchanged,
including all temporal and virtual endpoints. There are no duplicated
real vertices or internal edges in this construction, only two copies
of the artificial boundary. No edge is inferred from temporal coordinates.
\end{proof}
The hypothesis is necessary for this type of boundary-only resolution to
encode \emph{every} invisible chain by terminal parity: an internal odd
cycle has no terminal incidence under any such rewiring. The triangle
counterexample above supplies this obstruction. A more general certified
surface cut can duplicate internal vertices or subdivided edges, but must
prove both cost-nonincreasing lifts and projections of admissible chains;
no such local geometric construction is established here for all Lee
memory schedules. Use Theorem~\ref{thm:cl-cover} whenever the balance gate
fails. This is an efficient geometry-independent memory algorithm, with
its stated multi-source-run cost, not an unproved one-Dijkstra claim.

\begin{corollary}[Representative invariance]\label{cor:cl-gauge}
Replacing $L_2$ by $L_2+g^{\mathsf T}D$ while holding each residual cost
fixed leaves $\phi_c^{\rm circ}$ unchanged. Any two valid cut constructions
for equivalent cochains compute that same value. The covers are isomorphic
by $(v,a)\mapsto(v,a+g(v))$ with $g(b_*)=0$.
\end{corollary}
\begin{proof}
The extra term vanishes on every feasible difference. The displayed
vertex permutation carries each lifted edge to its gauged lift and keeps
its weight. Apply the two exact-minimum theorems.
\end{proof}
This invariance concerns the same original coverage. Re-running an
arbitrary growth algorithm on different cut graphs can change coverage
and is not an instance of this result.

\section{Cluster contraction with preserved logical information}
\label{sec:cl-clusters}
\begin{theorem}[Transport of interval coverage]\label{thm:cl-transport}
Choose finite nonnegative radii at original vertices of $G_c^\bullet$
and form the union of metric balls in its interval realization.
For $e=uv$ let
\[
 h_v=\max(0,\max_s(r_s-d(s,v))),\qquad
 \bar\omega_c(e)=\max(0,w_e-h_u-h_v).
\]
This is exactly the uncovered length of the original edge. A loop uses
both endpoints at the same vertex and hence $\max(0,w_e-2h_u)$.
Assign each cover copy the residual cost of its source edge. In a cut,
assign unchanged edges their source residual cost; if an edge is subdivided,
assign its pieces their actual uncovered interval lengths, whose sum is
the source residual. With logical labels and source IDs retained, these
operations preserve the minimum in \eqref{eq:cl-swim}.
\end{theorem}
\begin{proof}
Every route to the interior of an edge enters through one endpoint.
Thus the covered prefixes and suffixes have lengths $h_u,h_v$, clipped
to the edge interval; their union proves the formula, also for loops.
Subdivide at coverage endpoints, give covered pieces zero cost, and keep
a subdivision label assignment whose XOR equals $\lambda_e$.
A cycle in the subdivision uses either all or none of an original edge
because new vertices have degree two. The cover's projection/lifting
proof therefore applies with exactly the same costs and labels.
Equivalently, lift the entire covered subset to both sheets, then contract
its connected components \emph{in the cover}, retaining expansion paths.
A route through a contracted component expands at zero cost on that sheet
structure. This proves the residual transport without new decoder growth.
For a valid cut use its stated chain correspondence with these same costs.
\end{proof}
Spatial, timelike and diagonal edges have identical treatment. Virtual
vertices may be metric centers if they carry the stage-2 syndrome.
Half-edges are actual intervals with their own labels, even when their
end-time coordinates coincide. Clusters touching boundaries do not create
extra positive coverage by themselves. All coverage must be computed
from the explicitly chosen original metric, not from merged component IDs.

There is an essential qualification to the word contraction: a covered
component can contain an odd cycle. Erasing that component's labels in
the base graph loses a zero-cost logical witness. One may retain all
original zero-cost labelled edges (as here), or contract in the cover with
expansion data. A single unlabelled quotient vertex is insufficient.
For example two fully covered parallel half-edges with different labels
have $\phi=0$, although collapsing them into one unlabelled edge loses it.
A zero answer is returned with a nonempty mod-two logical witness.

If radii are certified boundary-normalized odd-cut radii, $r_s\le d(s,b_*)$.
Such a ball cannot propagate a positive distance through the artificial
boundary. Splitting that boundary as in Theorem~\ref{thm:cl-cut} therefore
preserves covered portions, as in Proposition~\ref{prop:dual}; components
must still be recomputed on the cut. For arbitrary production radii this
commutation is not automatic. Define coverage first on $G_c^\bullet$ and
transport it; or record a different, fully specified coverage convention.
In the controlled Phase-1 reduction below the convention must match
Phase~1 exactly. Missing radii mean an unavailable cluster score, not zero.

Meister's Fig.~2 caption explicitly allows a three-dimensional decoding
graph with faulty measurements \cite[p.~6]{Meister}. The interval-ball
construction and odd-cut reasoning are graph-theoretic adaptations;
the logical quotient must still be supplied separately. Neither UF growth
nor the exported Sparse-Blossom final-defect balls are thereby certified.

\section{Certified circuit-level representative bound}
\label{sec:cl-gap}
For this theorem suppose the completed graph is connected after removal
of inactive isolated rows (loops are allowed), or apply the construction
componentwise, using ordinary even-syndrome $T$-joins in components
without a boundary. For clarity the proof uses a single connected graph
with boundary. Let $\Sigma=\operatorname{supp}s$. On its syndrome metric
closure plus $b_*$, take the nonnegative boundary-normalized odd-cut dual
of Definition~\ref{def:dual}, now with distances in $G_c^\bullet$.
Thus $y_S\ge0$ for odd $S\subseteq\Sigma$,
\[
 l_{uv}(y)\le d(u,v),\quad r_u=\sum_{S\ni u}y_S\le d(u,b_*),
 \quad Y=\sum_Sy_S=w(f),\quad Df=s.
\]
The equality is required in exact arithmetic for the exact supplied
weights. The parity-polyhedron reduction of Proposition~\ref{prop:dual}
gives such a certificate for any exact minimum $f$. Nonnegative loops
can be removed from a minimum; selected zero-cost loops do not affect it.
The explicit companion LP remains exponentially represented in syndrome
size, separate from the polynomial cover postprocessing.

\begin{theorem}[Fixed-fiber certified gap]\label{thm:cl-gap}
With the specified optimal certificate and its interval coverage, set
\[
 W_{c,\mathrm{base}}^{\rm circ}=w(f),\qquad
 W_{c,\mathrm{opp}}^{\rm circ}
 =\min_{Dm=s,\ L_2(m+f)=1}w(m).
\]
If the opposite class exists, then
\begin{equation}
 W_{c,\mathrm{opp}}^{\rm circ}-W_{c,\mathrm{base}}^{\rm circ}
 \ge\phi_c^{\rm circ}.                              \label{eq:cl-gap}
\end{equation}
No two-terminal cut or spacetime embedding is a hypothesis.
\end{theorem}
\begin{proof}
Pair half-edges of a feasible $m$ at each real vertex, leaving exactly
one unpaired at each syndrome vertex and all boundary half-edges
unpaired. This decomposes $m$ into edge-disjoint trails and closed trails,
including loops. Each defect occurs as a trail endpoint once. A trail
between defects $u,v$ has covered length at least
$\min(\text{length},r_u+r_v)\ge l_{uv}(y)$; a defect-boundary trail has
covered length at least $r_u$. For every odd $S$, at least one endpoint
pair exits $S$, so summing these inequalities counts $y_S$ at least once.
Other trails contribute nonnegative length. Hence covered weight is at
least $Y$ for every feasible $m$, exactly as in Lemma~\ref{lem:coverage}.
For $f$, $Y=w(f)$ forces $\bar\omega_c(f)=0$. Therefore
\[
 w(m)-w(f)\ge\bar\omega_c(m)
 =\bar\omega_c(m+f)\ge\phi_c^{\rm circ}
\]
for every opposite candidate, since $D(m+f)=0$ and $L_2(m+f)=1$.
Take the minimum. The argument depends on the DEM logical map, not on
physical face generation or planar geometry.
\end{proof}
This is an exact minimum-representative cost inequality in a fixed
stage~1 fiber. It is not a summed-class likelihood ratio, a full
comparative-decoder gap, or a calibration theorem. Current exported
radii use the named final-defect metric-ball convention and do not export
$y_S$, its original defect memberships, verified feasibility, or exact
primal/dual equality. Consequently \eqref{eq:cl-gap} is \emph{not certified}
for those production outputs. With a suboptimal zero dual, a three-edge
logical path can have residual cost three while a one-edge baseline and
two-edge opposite correction have gap one; the Phase-1 counterexample
already disproves that stronger claim.

\section{Exact perfect-measurement reduction}
\label{sec:cl-limit}
\begin{theorem}[Controlled code-capacity limit]\label{thm:cl-limit}
Prepare the standard memory ideally, insert one layer of independent
$X_q$ mechanisms with probabilities $0<p_q<1/2$ immediately before the
first ideal syndrome extraction, and make every subsequent operation,
measurement and readout noiseless. Retain only this fault layer and the
measured-sector observable. Identify redundant zero detector rows as
padding. Under inherited weights and identical Phase-1 coverage,
\[
 D_c^{\rm circ}=D_c,\qquad L_c^{\rm circ}|_{K_c}=\beta_c,
 \qquad \phi_c^{\rm circ}=\phi_c^{\rm Phase1},
 \label{eq:cl-limit}
\]
where the first equality uses the explicit qubit and row bijections below.
Certified normalized balls grown in the merged graph give the same
coverage as their Phase-1 resolved realization. Arbitrary supplied
Phase-1 coverage can instead be transported edgewise by definition.
\end{theorem}
\begin{proof}
An $X_q$ inserted before the first ideal measurement changes every later
$Z$-check outcome containing $q$ by the same bit. Consecutive parities
therefore cancel after the first detector layer, and final ancilla/data
consistency also cancels. The only nonzero detector column is $Hq$ on
first-layer $Z$ faces. Its final logical label is $\ell_Z^{\mathsf T}q$.
There are no other admissible mechanisms; setting their probabilities
to zero means removing them, not retaining finite-cost columns.

By the Phase-1 local incidence proof, the non-$c$ detector target set
of $q$ identifies its primal $c$ edge when one exists. Thus restricted
compression groups its two endpoint mechanisms, with probability
$p_u+p_v-2p_up_v$, and creates exactly the virtual row for that edge.
At the opposite corner the target set is empty and no virtual row is
created. The two endpoints have distinct $c$ incidences as in
Proposition~\ref{prop:cl-memory}; different physical $c$ edges have
distinct restricted incidences in the specified hexagonal patch.
All qubits pass the filter, and complete stage-2 targets remain distinct.
Consequently there is one retained stage-2 edge per qubit, with its own
$p_q$, and the row/column identification gives
$D_c^{\rm circ}=(H_c,M_c)T_c=D_c$.
On its kernel, $L_c^{\rm circ}z=\ell_Z^{\mathsf T}T_cz=\beta_c(z)$ by
Lemma~\ref{lem:sideparity}. Hence the two fundamental minima have
identical feasible supports and identical residual cost on each edge.
Theorem~\ref{thm:swim} proves the required equality with the Phase-1
terminal distance. For certified balls, the radius-to-boundary bound
prevents positive propagation across the merged boundary, and
Proposition~\ref{prop:dual} proves the covered-interval equality.
\end{proof}
This is a theorem for a specified limit, including any number of additional
noiseless rounds. Repeated noisy data layers are a different mechanism
space. Homotopy equivalence alone would not prove metric equality, and
independently chosen optimal duals need not give equal scores. The
reduction concerns a fixed stage~1 fiber; it does not require the
restricted weighted hard decoder to select the same fiber as an unweighted
Phase-1 presentation.

\section{Validation contract and limits of acceptance}
\label{sec:cl-validation}
The following are implementation acceptance tests, not substitute proofs
and not a numerical study performed by this theory task.
\begin{enumerate}
\item \textbf{Exact reduction.} Construct the controlled single-layer limit
at $d=3,5,7$ in each color, compare $D,L_2$ under explicit row/edge maps,
and compare residual minima with identical coverage. Do not merely compare
sampled hard failure rates.
\item \textbf{Exhaustive small DEM.} For tiny retained memory DEMs at small
$d,T$, enumerate all mechanism subsets when feasible, evaluate $Dz,L_2z$
and residual cost, and compare with the cover. Save actual column counts;
if $2^{|E|}$ is too large, use explicitly identified small submodels or an
exact binary optimizer, without calling that full enumeration.
\item \textbf{Witness.} For every finite returned value, including zero,
check $Dz=0$, $L_2z=1$ and cost equality using original mechanism IDs.
For an empty logical class require infinity and no witness. Test isolated
odd loops, parallel labels, internal odd cycles and disconnected graphs.
\item \textbf{Stim.} With zero radii, positive-probability separator-free
graphlike mechanisms and unit weights, compare with
\texttt{DetectorErrorModel.shortest\_graphlike\_error}. This API minimizes
mechanism count, ignoring nonzero probabilities, and can treat separator
components as graphlike pieces \cite{StimAPI}. It is not an oracle for
nonuniform log odds or for correlated component costs. Assert witness
parity, not identity of tied minima. The API's prose ``less than two''
symptoms conflicts with its two-detector examples; the operational
criterion here is at most two detectors.
\item \textbf{Cut and representative checks.} Verify the internal balance
condition before using a two-terminal cut. Compare it with the cover;
apply arbitrary detector-row gauges and compare again with fixed residuals.
Reject cuts lacking a parity-preserving witness projection.
\item \textbf{Temporal audit.} Retain first/final record parities, source
faults and observable labels. Test temporal half-edges with both labels,
and reject time-only or endpoint-role-only logical classifications.
\item \textbf{Compression and provenance.} Compare full-target XOR maps;
include duplicate targets, incompatible labels on parallel endpoints,
empty-detector logical mechanisms, filtering losses and mixed-sector
fallback collisions. Check manager-to-stage-2 logical and detector
identities for each retained column.
\item \textbf{Contraction.} Compare interval subdivision with transported
residual edges, including partial coverage, loops, overlapping balls,
boundary contacts and covered odd cycles. Never erase their logical labels.
\item \textbf{Hard invariance.} Freeze input matrices, probabilities,
ordering, merge policy and syndrome arrays. Compare hard predictions,
weights, selected color and failure labels before/after analysis on the
same records. Circuit SO is not enabled until this future gate passes.
\item \textbf{Certificate.} If exact dual data are supplied, verify all
odd-cut constraints and equality to the actual labelled-graph optimum;
then compare exhaustive representative gaps. Without that certificate
report the theorem check as inapplicable, not passed.
\end{enumerate}
Remaining open issues are the local prism/fault-complex lift for the actual
Lee compressed graph, a universally valid geometric one-search cut for
all schedules, and production optimal-dual export. The DEM quotient,
cover, gated cut, coverage transport, controlled reduction and conditional
representative bound do not depend on resolving those stronger claims.
No aggregation theorem, likelihood calibration, sliding-window rule,
production decoder change, or performance benchmark is included.

\section{Executable DEM-to-decoder construction specification}
\label{sec:cl-algorithm}
This construction specifies a read-only soft-output adapter; the compact
implementation interface is given in Section~\ref{sec:cl-implementation}.
The existing hard path is authoritative, including its effective-model
choices. For a new circuit require explicit detector/observable metadata;
coordinates alone are insufficient to infer Pauli type, color or logical
boundary. The following records and gates are part of the algorithm.
\begin{enumerate}
\item \textbf{Freeze the experiment.} Inputs: Stim circuit, memory basis,
$d,T$, schedule, noise parameters and conversion options. Outputs: circuit
hash, package SHAs, original DEM (when exact extraction is supported),
manager's separated circuit/DEM, and an approximation ledger. Do not replace
the physical sampling circuit with the separated decoding circuit.
\item \textbf{Index targets and sources.} Flatten repeats and detector
offsets. Save mechanism IDs as (model hash, flattened instruction index),
probability, full detector/observable parity sets, separator-group IDs,
and available circuit error-location explanations. Save detector ID,
record parity support, full coordinates, time, check Pauli and color;
save observable ID, measured Pauli and record parity support. Keep absent
geometric locations explicitly absent. Validate all array dimensions.
\item \textbf{Materialize the actual decomposition.} Read the frozen
manager's $H_1,H_2,p_1,p_2$, stage-2 observable matrix, sort permutation,
and error-map matrices. Tag every virtual row by restricted edge ID and
its entire source target set $N$. Distinguish full-target compression
classes from selected representative maps. Record filtered sources and
all probability changes. For a paper-algorithm reference, use
\eqref{eq:cl-compression}--\eqref{eq:cl-decomp} and keep it separately
identified; it must not replace the hard path implicitly.
\item \textbf{Run unchanged stage~1.} Supply the historical restricted
check matrix, $\log((1-p_1)/p_1)$ and its same row-masked measured syndrome
to PyMatching. Its predicted edge vector $t$ indexes the virtual rows.
Check restricted syndrome feasibility, with original row IDs restored.
Keep the exact column ordering and hard merge/tie policy.
\item \textbf{Build the same stage~2 input.} Mask physical detector bits
to color $c$ and append $t$, preserving the historical padding. Supply
unchanged $H_2$, log-odds weights and syndrome to the existing stage-2
PyMatching call. Default check-matrix fault IDs represent stage-2 columns;
physical observable labels remain in the external $L_2$ matrix and need
not be added as hard parity constraints. Keep the ordinary prediction
$f$ and weight separate from analysis.
\item \textbf{Build the labelled analysis model.} From each $H_2$ column
retain its ID, complete nonzero endpoint rows, source effective DEM IDs,
probability, exact chosen weight, $L_2$ label, and physical/virtual endpoint
metadata. Audit $D f=s$ and $L_2f$ against the unchanged outer frame map.
Remove only zero padding rows. Keep parallel mechanisms and zero-detector
logical loops even if hard graph construction omits or merges them. Check
$L_2\notin\operatorname{rowspan}D$ or construct the equivalent odd-cycle
witness. If this fails, return ``no opposite class'' and infinity.
\item \textbf{Audit hard-graph compatibility.} PyMatching's internal parallel
merge and weight quantization can differ from the retained multigraph.
Record the map and recompute $w(f)$ from retained columns. The geometric
metric can still be defined on the retained graph with an explicitly
chosen ball convention, but the bound flag requires an optimal $f$ and
certificate for \emph{that same} graph and weights. An internal hard weight
alone is not this audit.
\item \textbf{Obtain growth.} Read existing growth data without altering
hard decoding, with original defect IDs, radii, units, normalization and
the named convention. If unavailable return ``growth unavailable''.
A separately authorized exact certificate extractor/companion solver
would additionally provide $y_S$, feasibility and equality to $w(f)$;
current final-defect radii do not. Set the certification flag accordingly.
\item \textbf{Compute residuals once.} Complete half-edges at $b_*$ and
zero-detector columns as loops. For the declared original ball metric,
compute $h_v$ by a temporary-source Dijkstra with initial costs
$R-r_v$, $R=\max r_v$. The temporary source is discarded. Evaluate every
$\bar\omega_c(e)$, retaining original IDs. A supplied interval-coverage
convention instead unions intervals per source edge explicitly.
\item \textbf{Choose the proven logical analysis.} Test balance of the
internal graph. On success the cut of Theorem~\ref{thm:cl-cut} uses one
terminal search. Otherwise build the cover and minimize same-base-vertex
opposite-sheet distances as in Theorem~\ref{thm:cl-cover}. All copies inherit
their source residual; no new matching or decoder growth is performed
on the cover/cut. Their vertices, sheets, boundaries and temporary search
arcs are analysis-only and never passed as new hard-decoder checks.
\item \textbf{Return and verify.} Return $\phi_c^{\rm circ}$, original-edge
mod-two witness, graph/weight/growth identifiers, class-existence status
and bound-certification flag. Verify witness incidence, frame parity and
cost independently. Return ordinary correction, ordinary weight and
existing branch selection unchanged. Report per-color fixed-fiber scope;
exclude open future boundaries and sliding-window decoding. The separately
authorized closed-memory implementation does not change these theorem hypotheses.
\end{enumerate}
